\section{Discussion}
\label{sec:discussion}

\paragraph{Why clones, not gloss.}
The most instructive result is the family asymmetry in the forced-choice
probe. Search-optimized pages and review-flooded pages are rejected at high
rates: agents ground their choices in specification text, and when the
specifications do not support the instruction, marketing surface does not
save the page. Specification clones invert this. A clone inherits the gold
listing's structure, alters one or two critical attributes, and undercuts the
price; at decision time it is a coin flip against the true product. This
recasts the threat model: the dangerous form of ambient synthetic content is
not the attention-grabbing page but the \emph{near-invariant forgery}, which
exploits exactly the behavior --- specification grounding --- that makes
agents robust to gloss. It also explains why detector-style defenses based on
surface statistics can only partially help: the clone's text statistics look
organic because its text \emph{is} organic, copied from a real product.

\paragraph{The economics of silent erosion.}
Success rates stayed nearly flat while regret climbed and confidence rose.
An operator watching task success --- the metric most dashboards show ---
would see nothing. The losses are absorbed inside ``successful'' sessions as
attribute mismatches and overpayment, which is precisely why an ambient
contamination problem can grow for years without surfacing as an incident.
The confidence inflation compounds it: agents err more while trusting
themselves more, so downstream automation (reordering, auto-purchase budgets)
inherits correlated, concealed losses. We argue this pattern --- flat
success, rising regret, rising confidence --- should become the standard
diagnostic triplet for evaluating agents on live-content distributions, in
the same way calibration curves preceded accuracy as the trustworthy signal
for classifiers under distribution shift.

\paragraph{Provenance as an interface contract.}
Our defense results argue for a specific division of labor. Content-level
detection asks every consumer to re-litigate every page's authenticity, and
it fails exactly on the content that matters (near-copies of real pages).
Provenance-weighted ranking instead asks the supply side for one bit ---
verifiable origin --- and lets the agent treat its absence as a prior, not a
verdict. Because the bit is cheap to spoof in the model but expensive to
spoof in the world (a signature over seller identity and listing history),
it converts an unbounded text-classification arms race into a bounded
key-management problem. The noise we inject (15\% missing on real items,
15\% present on synthetic ones) is an intentionally pessimistic model of
real-world deployment, and recovery remains high; the remaining gap to the
oracle is the price of noisy provenance, not of the idea.

\paragraph{Limitations.}
Our agents are 7--8B open-weight models with a scripted action interface;
frontier models with stronger browsing may resist clones better --- though
the clone attack works \emph{because} of specification grounding, which
stronger models exhibit more strongly, so the vulnerability may persist or
worsen. The contamination families are generated by models of the same size
class as the agents; industrial content farms use larger generators, which
could be more persuasive, but our plausibility floor and manual audits
suggest the families are representative of real farm output. WebShop
instructions are variant-level and sometimes weakly aligned with the base
listing text, an inherent property of the benchmark that we handle with the
official best-variant scoring; this makes absolute success numbers
benchmark-specific but does not interact with the contamination manipulation,
which is randomized within task. Our dose--response uses a single item pool
per task; we mitigate pool-specific artifacts by keeping pools fixed across
$\varepsilon$ (paired design) and by releasing all pools and seeds. Finally,
provenance availability in the wild is an empirical question our 85/15 noise
model approximates crudely; measuring real provenance coverage on major
marketplaces is a natural next step.

\paragraph{Outlook.}
If agentic browsing becomes a dominant traffic source, the incentive to
optimize \emph{for agents} will exceed the incentive to optimize for human
attention, because agents can be measured and A/B-tested at scale. The
question is not whether the Web's supply side adapts to machine readers ---
it already is --- but whether machine readers will be able to tell what
adapted. Ambient contamination is the benign end of that adaptation spectrum;
deliberate manipulation is the malign end, and both flow through the same
missing infrastructure: content whose origin a buyer can verify.
